\section*{Solutions to Homework 3}

\begin{lstlisting}
$ bash codeHints/L03/setup_hw3.sh
$ cd practice/hw3
$ git log --oneline
e672c50 Update notse
6e1cc4f Buy fireworks
809893e Add recipe notes
2217470 Add shopping list
$ git restore list.txt
$ git restore --staged notes.txt
$ git commit --amend -m "Update notes"
$ git revert --no-edit 6e1cc4f
$ git status --short
 M notes.txt
$ printf 'fireworks=6e1cc4f\ncommits=4\nno-undo=restore\n' > ../hw3_answers.txt
\end{lstlisting}

(Your IDs for the first three commits match these because the setup script fixes author and dates. The amended and revert commits get new IDs.)

\textbf{3.1.} \code{git log --oneline} shows the ID \code{6e1cc4f} next to \code{Buy fireworks} and four lines, so \code{commits=4}. The answer file lives in \code{practice/}, outside the repository, so it does not show up as an untracked file.

\textbf{3.2.} \code{git diff list.txt} first shows the junk line \code{+oops delete me}. \code{git restore list.txt} removes it. This is the destructive command: it is correct here, because the line really was junk. You must do this \emph{before} the revert, because the revert changes \code{list.txt} too and Git will not overwrite your uncommitted edit.

\textbf{3.3.} \code{git restore --staged notes.txt} moves the \code{TODO: add syrup} line from ``staged'' back to ``modified''. The line stays in the file. Check with \code{git status --short}: the first column is now empty.

\textbf{3.4.} \code{git commit --amend -m "Update notes"} rewrites the last commit with the corrected message. The order matters: if you amend while \code{notes.txt} is still staged, the TODO line is pulled into the commit. That is why 3.3 comes first. If you made that mistake, run \code{setup\_hw3.sh --reset} and start again.

\textbf{3.5.} \code{git revert --no-edit 6e1cc4f} creates \code{Revert "Buy fireworks"} on top of the amended commit. The history has five commits now. It works without conflict because the commits after \code{Buy fireworks} only touched \code{notes.txt}.

\textbf{3.6.} The answer is \code{restore}. It discards uncommitted work for good. (\code{restore --staged} keeps your edit; \code{amend} leaves the old commit abandoned but recoverable, as Lecture 10 shows; and \code{revert} just adds a commit. None of the three deletes work from your files.)

\Needspace{8\baselineskip}
\textbf{Common mistakes.}
\begin{itemize}
  \item Doing 3.4 before 3.3: the TODO line gets committed. Reset and redo in the right order.
  \item Running \code{git revert} without \code{--no-edit} and getting stuck in the editor: press \code{Esc}, type \code{:wq}, Enter to accept the default message.
  \item Reverting the wrong commit (the typo commit instead of \code{Buy fireworks}). Look at the ID in \code{git log --oneline} before you revert.
  \item Using \code{git reset}, which this lecture did not teach: it can rewrite history in ways that this lecture does not cover, and the checker may flag the result.
\end{itemize}
